% gz-p10-04-2: 闭区间最值——y = x^3 - 3x 在 [-2, 3] 的 4 个候选点
\documentclass[tikz,border=4pt]{standalone}
\usepackage{ctex}
\usepackage{amsmath}
\usepackage{pgfplots}
\pgfplotsset{compat=1.18}

\begin{document}
\begin{tikzpicture}
\begin{axis}[
  width=12cm, height=9cm,
  axis lines=center,
  xlabel={$x$}, ylabel={$y$},
  every axis x label/.style={at={(axis cs:3.5,0)}, anchor=west},
  every axis y label/.style={at={(axis cs:0,20)}, anchor=south},
  xmin=-2.8, xmax=3.5,
  ymin=-3.5, ymax=20,
  xtick={-2,-1,1,3},
  ytick={-2,2,18},
  tick style={color=black},
  ticklabel style={font=\small},
  clip=true,
]
  % 曲线 y = x^3 - 3x，在 [-2, 3]
  \addplot[black, thick, domain=-2:3, samples=120] {x^3 - 3*x};
  \node[right, black] at (axis cs:2.5, 10) {$y=x^3-3x$};

  % 四个候选点
  % x = -2 端点: (-2, -2)
  \addplot[only marks, mark=*, mark size=2.8pt, blue] coordinates {(-2,-2)};
  \node[below left, blue, font=\footnotesize] at (axis cs:-2,-2) {端点$(-2,-2)$};

  % x = -1 极大值: (-1, 2)
  \addplot[only marks, mark=*, mark size=2.8pt, red] coordinates {(-1,2)};
  \node[above left, red, font=\footnotesize] at (axis cs:-1,2) {极大$(-1,2)$};

  % x = 1 极小值: (1, -2)
  \addplot[only marks, mark=*, mark size=2.8pt, blue] coordinates {(1,-2)};
  \node[below, blue, font=\footnotesize] at (axis cs:1,-2) {极小$(1,-2)$};

  % x = 3 端点: (3, 18) ← 最大值
  \addplot[only marks, mark=*, mark size=3pt, red] coordinates {(3,18)};
  \node[left, red, font=\footnotesize] at (axis cs:3,18) {端点$(3,18)$};
  \node[right, red, font=\small] at (axis cs:3,17) {\textbf{最大}};

  % 虚线
  \draw[gray, dashed] (axis cs:-2,-2) -- (axis cs:-2,0);
  \draw[gray, dashed] (axis cs:3,18) -- (axis cs:3,0);

  % 文字注释
  \node[font=\footnotesize, align=center] at (axis cs:0.5, 14)
    {比较 4 个候选点的 $y$\\取最大、最小};
\end{axis}
\end{tikzpicture}
\end{document}
